\section{State probabilities}

Using formula (\ref{stateFormula}), we can obtain the individual state probabilities. The first, $p_0(t)$ will differ between the catastrophe and killing cases: 
$$G(0 , t) = \begin{cases}
\frac{ab(1-\ee^{\beta(a-b) t})  } {  (b-a\ee^{\beta(a-b) t})   }, & \text{(catastrophe)}\\
1 - I(t) + \frac{ab(1-\ee^{\beta(a-b) t})  } {  (b-a\ee^{\beta(a-b) t})   }, & \text{(killing).}
\end{cases}$$
As stated above, the fraction term in each case represents the probability of reaching 0 by death, and the $1-I(t)$ pertains to reaching 0 by collapse in ``killing''.\par 
From here, the state probabilities, $p_1(t)$, $p_2(t)$, $p_3(t)\dots$, do not differ between the cases. Repeated differentiation of $G(z,t)$ with respect to $z$ reveals the formula,
$$\frac{\partial^n G}{\partial z^n} =n! \ee^{\beta(a-b)t}(a-b)^2  \frac{\lrb{1-\ee^{\beta(a-b)t}}^ {n-1} }{\lrb{b - a\ee^{\beta(a-b)t} - \lrb{1-\ee^{\beta(a-b)t}}z} ^ {n+1}} $$
Hence, setting $z=0$, and dividing by $n!$, the state probabilities for $n\geq1$ are described by,
\begin{equation}
\label{stateprob}
p_n(t) =\lrb{ \frac{ a-b }{{b - a\ee^{\beta(a-b)t}}}}^2 \ee^{\beta(a-b)t} \lrb{\frac{{1-\ee^{\beta(a-b)t}}}{{b - a\ee^{\beta(a-b)t}} } }^{n-1}.
\end{equation}
